\documentclass[10pt,a4paper]{amsart}
\usepackage[margin=19mm]{geometry}
\usepackage{amsmath,amssymb,mathtools}
\usepackage[scr=rsfs,cal=euler]{mathalfa}
\usepackage{xfrac}
\usepackage[unicode,hidelinks,bookmarks=false]{hyperref}
\newcommand{\ie}{i.e.\ }
\newcommand{\cf}{cf.\ }
\newcommand{\resp}{respectively\ }
\newcommand{\Subsection}{\subsection}
\providecommand{\nobreakdash}{\nobreak\hbox{-}}
\setlength{\emergencystretch}{2em}
\multlinegap0pt
\usepackage{bm}
\usepackage{bbm}
\usepackage{dsfont}
\usepackage{xspace}
\usepackage{cancel}
\usepackage{mathtools}
\usepackage{amsthm}
\makeatletter
\@ifundefined{proposition}{\newtheorem{proposition}{Proposition}}{}
\@ifundefined{theorem}{\newtheorem{theorem}{Theorem}}{}
\makeatother
\DeclareMathOperator*{\argmin}{arg\,min} 
\newcommand{\N}{\mathbb{N}}
\title[Cluster-Aware Matching via Laplacian Optimal Transport]{Cluster-Aware Matching via Laplacian Optimal Transport}
\author{Gabriel Samberg, YoonHaeng Hur, Yuehaw Khoo, Nir Sharon}
\date{Source: arXiv:2607.16178v1 (CC BY 4.0)}
\begin{document}
\maketitle

\noindent\emph{Source and licence.} Cluster-Aware Matching via Laplacian Optimal Transport. Version arXiv:2607.16178v1 (CC BY 4.0), distributed under the Creative Commons Attribution 4.0 International licence, as recorded in the arXiv licence metadata for this exact version and shown on the abstract page https://arxiv.org/abs/2607.16178v1. Full rights notice: licenses/arxiv-2607.16178-CC-BY-4.0.txt.\par\smallskip

\noindent\textit{Editorial scope.} Counted unit: the Proposition whose source label is \texttt{None} in the article (source0.tex lines 381--397, source label \texttt{None}), delivered together with the article's own complete original proof of it (source0.tex lines 398--450, one \texttt{proof} environment of 53 lines). No mathematics is written, repaired or re-derived: statement and proof are the article's own text line for line. Required context retained uncounted: thm:optimal\_coupling\_projection\_bound (lines 294--306); eq:optimal\_coupling\_projection\_bound (lines 302--305). Every definition, display and label of those lines is retained unchanged. Omitted from this package and not redistributed: 1--293, 306--380, 451--921. Nothing inside a retained excerpt is omitted or altered: every retained body is delivered exactly as the article's lines are written, so no omission receipt of any kind accompanies this package. Citations: the retained excerpts carry no citation command at all, so no bibliography entry of the article is transcribed and no reference is redistributed. Reference adaptation: none was needed and none was made; every reference inside the delivered unit resolves inside it, so no unresolved reference or citation is produced. Printed slips kept literal: no printed word, symbol, hypothesis or number of the counted statement or of its proof was altered. Layout and class: the article's own class is \texttt{article} with packages including geometry, hyperref, graphicx, subcaption, caption, amsmath, stmaryrd, amssymb, amsfonts, float, tikz, algorithm, algpseudocode, diagbox; the wrapper is the lane's pinned \texttt{amsart} editorial wrapper at 10pt on A4, which restates the theorem-like environments the retained text needs and adds the article's own macro definitions that the retained text needs (\texttt{\textbackslash N}, \texttt{\textbackslash argmin}), with extra wrapper packages (bm, bbm, dsfont, xspace, cancel, mathtools). The delivered numbering is the unit's own. Assets: no retained span contains a figure environment, a graphics-inclusion command, a PDF-page inclusion, a table or a TikZ picture; no image file is copied into this package, the package declares no asset dependency and no image file is redistributed with it. Licence: version arXiv:2607.16178v1 of this article is distributed under the Creative Commons Attribution 4.0 International licence, as recorded in the arXiv licence metadata for this exact version and shown on the abstract page https://arxiv.org/abs/2607.16178v1. A full rights notice is delivered at licenses/arxiv-2607.16178-CC-BY-4.0.txt. Characters: every retained excerpt is pure ASCII, so the delivered engine, XeLaTeX with its default Latin Modern text font, reports no missing character.
\begin{theorem}
    \label{thm:optimal_coupling_projection_bound}
    Let $\pi^\star$ be the solution to $\mathsf{LapOT}_{\lambda_x, \lambda_y, \lambda}(a, b, C, L_X, L_Y)$ with $\lambda_x, \lambda_y > 0$. Assume the following.
    \begin{itemize}
        \item[(i)] $L_X$ and $L_Y$ are induced by graphs with $r$ and $s$ connected components, where $\{A_\alpha\}_{\alpha = 1}^{r}$ and $\{B_\beta\}_{\beta = 1}^{s}$ are the corresponding partitions of the index sets $[n]$ and $[m]$, respectively.
        \item[(ii)] $a$ and $b$ are block-constant with respect to the partitions $\{A_\alpha\}_{\alpha = 1}^{r}$ and $\{B_\beta\}_{\beta = 1}^{s}$, respectively.
    \end{itemize}
    Let $L_X = \Phi_X \Lambda_X \Phi_X^\top$ and $L_Y = \Phi_Y \Lambda_Y \Phi_Y^\top$ be their orthogonal decompositions, where $\Lambda_X$ and $\Lambda_Y$ are diagonal matrices consisting of the eigenvalues $\mu_1^X \le \cdots \le \mu_n^X$ and $\mu_1^Y \le \cdots \le \mu_m^Y$ of $L_X$ and $L_Y$, respectively. For any $\ell \in [n]$ and $h \in [m]$, define the projection matrices as $P_{X, \ell} \triangleq \Phi_{X, \ell} \Phi_{X, \ell}^\top$ and $P_{Y, h} \triangleq \Phi_{Y, h} \Phi_{Y, h}^\top$, where $\Phi_{X, \ell}$ and $\Phi_{Y, h}$ are the submatrices consisting of the first $\ell$ and $h$ columns of $\Phi_X$ and $\Phi_Y$, respectively. Then, $P_{X, r} \pi^\star$ and $\pi^\star P_{Y, s}$ are obtained by block-averaging the rows and columns of $\pi^\star$ with respect to the partitions $\{A_\alpha\}_{\alpha = 1}^{r}$ and $\{B_\beta\}_{\beta = 1}^{s}$, respectively. Moreover, we have
    \begin{equation}
        \label{eq:optimal_coupling_projection_bound}
        \|\pi^\star - P_{X, r} \pi^\star\|_{\mathrm{F}}^2 \le \frac{\|P_{X, r} C - C\|_\infty}{\lambda_x \mu_{r + 1}^X} \quad \text{and} \quad \|\pi^\star - \pi^\star P_{Y, s}\|_{\mathrm{F}}^2 \le \frac{\|C P_{Y, s} - C\|_\infty}{\lambda_y \mu_{s + 1}^Y}.
    \end{equation}
\end{theorem}

\begin{proposition}
    In Theorem \ref{thm:optimal_coupling_projection_bound}, we deduce the following.
    \begin{itemize}
        \item[(i)] If $\lambda_x \to \infty$, then $\pi^\star$ converges as follows:
        \begin{equation*}
            \lim_{\lambda_x \to \infty} \pi^\star = \argmin_{\pi \in \Pi_{a, b}} \langle \pi, P_{X, r} C \rangle + \lambda_y \langle \pi, \pi L_Y \rangle - \lambda H(\pi).
        \end{equation*}
        \item[(ii)] If $\lambda_y \to \infty$, then $\pi^\star$ converges as follows:
        \begin{equation*}
            \lim_{\lambda_y \to \infty} \pi^\star = \argmin_{\pi \in \Pi_{a, b}} \langle \pi, C P_{Y, s} \rangle + \lambda_x \langle \pi, L_X \pi \rangle - \lambda H(\pi).
        \end{equation*}
        \item[(iii)] If $\lambda_x, \lambda_y \to \infty$, then $\pi^\star$ converges as follows:
        \begin{equation*}
            \lim_{\lambda_x, \lambda_y \to \infty} \pi^\star = \argmin_{\pi \in \Pi_{a, b}} \langle \pi, P_{X, r} C P_{Y, s} \rangle - \lambda H(\pi).
        \end{equation*}
    \end{itemize}
\end{proposition}

\begin{proof}
    To show (i), let $\pi^\star_{\lambda_x}$ denote the optimal coupling for a given $\lambda_x$. Now, consider a sequence $\{\lambda_x^{(k)}\}_{k \in \mathbb{N}}$ such that $\lambda_x^{(k)} \to \infty$ as $k \to \infty$. Since $\Pi_{a, b}$ is compact, there exists a convergent subsequence of $\{\pi^\star_{\lambda_x^{(k)}}\}_{k \in \mathbb{N}}$; by swapping the original sequence with its convergent subsequence, we may assume that $\{\pi^\star_{\lambda_x^{(k)}}\}_{k \in \mathbb{N}}$ converges to some $\tilde{\pi} \in \Pi_{a, b}$. We now claim 
    \begin{equation*}
        \tilde{\pi} = \argmin_{\pi \in \Pi_{a, b}} \langle \pi, P_{X, r} C \rangle + \lambda_y \langle \pi, \pi L_Y \rangle - \lambda H(\pi) =: \pi^\star_\infty.
    \end{equation*}
    We first note that $\pi^\star_\infty = P_{X, r} \pi^\star_\infty$. To see this, as in the proof of Theorem \ref{thm:optimal_coupling_projection_bound}, observe that $P_{X, r} \pi^\star_\infty \in \Pi_{a, b}$, $H(P_{X, r} \pi^\star_\infty) \ge H(\pi^\star_\infty)$, and $\langle P_{X, r} \pi^\star_\infty, P_{X, r} \pi^\star_\infty L_Y \rangle \le \langle \pi^\star_\infty, \pi^\star_\infty L_Y \rangle$. Hence,
    \begin{equation*}
        \langle P_{X, r} \pi^\star_\infty, P_{X, r} C \rangle + \lambda_y \langle P_{X, r} \pi^\star_\infty, P_{X, r} \pi^\star_\infty L_Y \rangle - \lambda H(P_{X, r} \pi^\star_\infty) \le \langle \pi^\star_\infty, P_{X, r} C \rangle + \lambda_y \langle \pi^\star_\infty, \pi^\star_\infty L_Y \rangle - \lambda H(\pi^\star_\infty),
    \end{equation*}
    where we also use $\langle P_{X, r} \pi^\star_\infty, P_{X, r} C \rangle = \langle \pi^\star_\infty, P_{X, r} C \rangle$. By the uniqueness of the minimizer $\pi^\star_\infty$, we conclude that $\pi^\star_\infty = P_{X, r} \pi^\star_\infty$.


    For notational convenience, let $\pi^\star_{\lambda_x^{(k)}} = \pi^\star_k$ for $k \in \N$. Then, by the optimality of $\pi^\star_k$, we have
    \begin{equation}
        \label{eq:optimality_limit}
        \langle \pi^\star_k, C \rangle + \lambda_x^{(k)} \langle \pi^\star_k, L_X \pi^\star_k \rangle + \lambda_y \langle \pi^\star_k, \pi^\star_k L_Y \rangle - \lambda H(\pi^\star_k) \le \langle \pi^\star_\infty, C \rangle + \lambda_y \langle \pi^\star_\infty, \pi^\star_\infty L_Y \rangle - \lambda H(\pi^\star_\infty),
    \end{equation}
    where we use $\langle \pi^\star_\infty, L_X \pi^\star_\infty \rangle = 0$ from $\pi^\star_\infty = P_{X, r} \pi^\star_\infty$. Now, we claim 
    \begin{equation}
        \label{eq:limit_objective}
        \lim_{k \to \infty} \left(\langle \pi^\star_k, C \rangle + \lambda_x^{(k)} \langle \pi^\star_k, L_X \pi^\star_k \rangle + \lambda_y \langle \pi^\star_k, \pi^\star_k L_Y \rangle - \lambda H(\pi^\star_k)\right) = \langle \tilde{\pi}, P_{X, r} C \rangle + \lambda_y \langle \tilde{\pi}, \tilde{\pi} L_Y \rangle - \lambda H(\tilde{\pi}).
    \end{equation}
    To see this, use \eqref{eq:optimal_coupling_projection_bound} to deduce that 
    \begin{equation*}
        \lim_{k \to \infty} P_{X, r} \pi^\star_k = \lim_{k \to \infty} \pi^\star_k = \tilde{\pi}.
    \end{equation*}
    Hence, 
    \begin{equation}
        \label{eq:limit_temp1}
        \lim_{k \to \infty} \langle \pi^\star_k, C \rangle = \lim_{k \to \infty} \langle P_{X, r} \pi^\star_k, C \rangle = \lim_{k \to \infty} \langle \pi^\star_k, P_{X, r} C \rangle = \langle \tilde{\pi}, P_{X, r} C \rangle
    \end{equation}
    and
    \begin{equation}
        \label{eq:limit_temp2}
        \lim_{k \to \infty} \left(\lambda_y \langle \pi^\star_k, \pi^\star_k L_Y \rangle - \lambda H(\pi^\star_k)\right) = \lambda_y \langle \tilde{\pi}, \tilde{\pi} L_Y \rangle - \lambda H(\tilde{\pi}).
    \end{equation}
    Meanwhile, as in the proof of Theorem \ref{thm:optimal_coupling_projection_bound}, we have the following from the optimality of $\pi^\star_k$:
    \begin{equation*}
        \lambda_x^{(k)} \langle \pi^\star_k, L_X \pi^\star_k \rangle \le \langle P_{X, r} \pi^\star_k, C \rangle + \lambda_y \langle P_{X, r} \pi^\star_k, P_{X, r} \pi^\star_k L_Y \rangle - \lambda H(P_{X, r} \pi^\star_k) - \langle \pi^\star_k, C \rangle - \lambda_y \langle \pi^\star_k, \pi^\star_k L_Y \rangle + \lambda H(\pi^\star_k),
    \end{equation*}
    where the right-hand side converges to $0$ as $k \to \infty$ by the previous two limits. Hence, 
    \begin{equation}
        \label{eq:limit_temp3}
        \lim_{k \to \infty} \lambda_x^{(k)} \langle \pi^\star_k, L_X \pi^\star_k \rangle = 0.
    \end{equation}
    Combining \eqref{eq:limit_temp1}, \eqref{eq:limit_temp2}, and \eqref{eq:limit_temp3}, we obtain \eqref{eq:limit_objective}. Then, combining \eqref{eq:optimality_limit} and \eqref{eq:limit_objective}, we have
    \begin{equation*}
        \langle \tilde{\pi}, P_{X, r} C \rangle + \lambda_y \langle \tilde{\pi}, \tilde{\pi} L_Y \rangle - \lambda H(\tilde{\pi}) \le \langle \pi^\star_\infty, P_{X, r} C \rangle + \lambda_y \langle \pi^\star_\infty, \pi^\star_\infty L_Y \rangle - \lambda H(\pi^\star_\infty),
    \end{equation*}
    where we use $\langle \pi^\star_\infty, C \rangle = \langle \pi^\star_\infty, P_{X, r} C \rangle$ from $\pi^\star_\infty = P_{X, r} \pi^\star_\infty$. By the uniqueness of the minimizer $\pi^\star_\infty$, we conclude that $\tilde{\pi} = \pi^\star_\infty$.

    Hence, we have shown that for any sequence $\{\lambda_x^{(k)}\}_{k \in \mathbb{N}}$ such that $\lambda_x^{(k)} \to \infty$, the corresponding sequence of optimal couplings $\{\pi^\star_{\lambda_x^{(k)}}\}_{k \in \mathbb{N}}$ has a convergent subsequence that converges to $\pi^\star_\infty$. As the limit is independent of the choice of the sequence $\{\lambda_x^{(k)}\}_{k \in \mathbb{N}}$, we conclude that $\pi^\star_{\lambda_x} \to \pi^\star_\infty$ as $\lambda_x \to \infty$. This proves (i). The proofs of (ii) and (iii) follow similarly.
\end{proof}

\end{document}
